% ---------------------------------------------------------------------
%  02 Introduction
% ---------------------------------------------------------------------
\section{Introduction}
\label{sec:introduction}

\Todo{Opening paragraph: why 3D reconstruction of tennis scenes matters
(coaching, broadcast augmentation, officiating, biomechanics), and why it is
hard (small fast ball, motion blur, severe scale variation, occlusion,
uncalibrated broadcast cameras).}

\Todo{Second paragraph: the gap. Prior work treats ball tracking, court
registration, and player pose as separate problems, each evaluated in its own
2D image space. Nothing produces a single metric court-coordinate
reconstruction of ball and players jointly.}

\Todo{Third paragraph: our approach in one breath. We decompose the problem
into per-view 2D perception (ball, court keypoints, player pose) followed by
learned lifting modules that map 2D observations into court coordinates
(\BLCS{} for the ball, \PLCS{} for player root position and orientation), and
finally fuse \PLCS{} root transforms with \GVHMR{} body pose to place
articulated players in the court frame. \Cref{fig:teaser} illustrates the
output.}

% --- Teaser figure ----------------------------------------------------
\begin{figure}[t]
  \centering
  % \includegraphics[width=\linewidth]{assets/figures/teaser.pdf}
  \fbox{\parbox[c][0.28\textheight][c]{\linewidth}{\centering
    \Todo{teaser.pdf --- input multi-view frames (left) $\rightarrow$
    reconstructed 3D ball trajectory and player meshes in court coordinates
    (right)}}}
  \caption{\Todo{Teaser caption. State what the reader is looking at and what
  is novel about it, not how it was made.}}
  \label{fig:teaser}
\end{figure}

\Todo{Fourth paragraph: why this decomposition rather than end-to-end. The
court plane gives a metric anchor that is observable in every view; committing
to it early makes the lifting problem well posed and the failure modes
inspectable.}

\subsection*{Contributions}

\begin{itemize}
  \item \Todo{Contribution 1 --- the pipeline / system claim.}
  \item \Todo{Contribution 2 --- \BLCS{}: 2D ball + court $\rightarrow$ 3D ball
        trajectory, and what makes it work (temporal context, physics prior).}
  \item \Todo{Contribution 3 --- \PLCS{}: 2D pose + court $\rightarrow$ player
        root translation and rotation, including the rotation
        parameterisation that avoids $180^\circ$ ambiguity.}
  \item \Todo{Contribution 4 --- empirical study: datasets, ablations, and the
        sim-to-real transfer result.}
\end{itemize}
